\heading{经济学原理（II）-模拟题2-期中}
\noteinfo{题目和答案由 AI 依据样题与第 1--6 讲内容生成；本卷比模拟题 1 更强调综合性与“为什么错”的辨析。}

\instructions{建议答题时间 120 分钟。请特别注意：本卷多处同时考查“概念口径 + 图形直觉 + 代数计算”。}

\textbf{一、单项选择题}

\begin{question}
下列关于机会成本的表述，最准确的是
\begin{choiceoptions}[2]
  \optionitem{任何选择的机会成本都等于货币支出}
  \optionitem{机会成本等于放弃掉的所有选项价值之和}
  \optionitem{机会成本等于所放弃的最佳备选方案的价值}
  \optionitem{机会成本只存在于市场交易中}
\end{choiceoptions}
\begin{solution}
机会成本是为了得到某物而放弃的最佳备选项价值，不一定等于货币成本。
\anstext{C}
\end{solution}
\end{question}

\begin{question}
若某劳动者因为先接受了更高教育而获得更高工资，这一现象最直接体现的是
\begin{choiceoptions}[2]
  \optionitem{补偿性工资差异}
  \optionitem{效率工资}
  \optionitem{人力资本}
  \optionitem{道德风险}
\end{choiceoptions}
\begin{solution}
更高教育提升生产率并带来更高工资，体现的是人力资本理论。
\anstext{C}
\end{solution}
\end{question}

\begin{question}
在买方垄断的劳动力市场中，最低工资有时可能提高就业，原因在于
\begin{choiceoptions}[2]
  \optionitem{最低工资降低了工人的边际产品}
  \optionitem{最低工资可能使企业的边际雇佣成本下降到更接近工资本身}
  \optionitem{最低工资自动消除了失业}
  \optionitem{最低工资让劳动需求曲线右移}
\end{choiceoptions}
\begin{solution}
买方垄断下企业面临的边际雇佣成本高于工资；适度最低工资可能压低这一边际成本，从而提高就业。
\anstext{B}
\end{solution}
\end{question}

\begin{question}
下列哪一项会计入中国 2026 年 GDP，但\textbf{不会}计入 2026 年中国 GNP 的增加？
\begin{choiceoptions}[2]
  \optionitem{中国居民在国外获得的股息}
  \optionitem{驻华外资企业在中国境内创造的新增产出，其收益最终汇回国外母公司}
  \optionitem{中国政府向居民发放失业补助}
  \optionitem{中国居民购买二手房支付给卖家的房款}
\end{choiceoptions}
\begin{solution}
境内外资企业在中国境内生产计入 GDP，但其收益最终归外国居民所有，因此不增加中国 GNP。
\anstext{B}
\end{solution}
\end{question}

\begin{question}
关于拉氏指数、帕氏指数与费雪指数，下列说法正确的是
\begin{choiceoptions}[2]
  \optionitem{费雪指数完全不需要价格数据}
  \optionitem{拉氏指数通常使用当期数量作权重}
  \optionitem{帕氏指数通常使用当期数量作权重}
  \optionitem{三者在任何情况下都相等}
\end{choiceoptions}
\begin{solution}
拉氏指数使用基期数量，帕氏指数使用当期数量，费雪指数是二者的几何平均。
\anstext{C}
\end{solution}
\end{question}

\begin{question}
在索罗模型中，黄金规则资本存量对应的是
\begin{choiceoptions}[2]
  \optionitem{使稳态总产出最大}
  \optionitem{使稳态人均消费最大}
  \optionitem{使折旧为零}
  \optionitem{使储蓄率恒等于 1}
\end{choiceoptions}
\begin{solution}
黄金规则追求的是最大稳态人均消费，而不是最大产出。
\anstext{B}
\end{solution}
\end{question}

\begin{question}
若可贷资金市场中企业获得投资税收抵免，则最直接的影响是
\begin{choiceoptions}[2]
  \optionitem{储蓄供给右移}
  \optionitem{投资需求右移}
  \optionitem{储蓄供给左移}
  \optionitem{投资需求左移}
\end{choiceoptions}
\begin{solution}
投资税收抵免提高投资回报，增强企业投资意愿，因此投资需求右移。
\anstext{B}
\end{solution}
\end{question}

\begin{question}
若预期通胀率下降而名义利率不变，则真实利率会
\begin{choiceoptions}[2]
  \optionitem{下降}
  \optionitem{上升}
  \optionitem{保持不变}
  \optionitem{一定变成零}
\end{choiceoptions}
\begin{solution}
由 $r\approx i-\pi^e$，名义利率不变而预期通胀下降，则真实利率上升。
\anstext{B}
\end{solution}
\end{question}

\begin{question}
关于有效市场假说与资产泡沫，下列最恰当的是
\begin{choiceoptions}[2]
  \optionitem{只要市场有效，泡沫就绝不可能出现}
  \optionitem{有效市场假说说明价格一定等于基本面值}
  \optionitem{即使公开信息被迅速反映，价格仍可能因预期、行为偏差或制度摩擦阶段性偏离基本面}
  \optionitem{泡沫的存在自动否定一切有效市场观点}
\end{choiceoptions}
\begin{solution}
有效市场强调信息反映速度，不等于排除一切阶段性偏离。
\anstext{C}
\end{solution}
\end{question}

\begin{question}
保险市场中的道德风险是指
\begin{choiceoptions}[2]
  \optionitem{投保后被保险人变得更不谨慎}
  \optionitem{高风险人群更愿意购买保险}
  \optionitem{保险公司一定会亏损}
  \optionitem{所有风险都能被分散掉}
\end{choiceoptions}
\begin{solution}
高风险人群更愿意投保是逆向选择；投保后行为更不谨慎才是道德风险。
\anstext{A}
\end{solution}
\end{question}

\textbf{二、简答题}

\begin{question}
为什么“进口消费品价格上涨”可能使 CPI 上升得比 GDP 平减指数更快？

\begin{solution}
CPI 关心的是居民消费篮子的生活费用，只要居民消费中包含进口品，进口消费品涨价就会直接推高 CPI。GDP 平减指数关心的是本国生产的最终产品和服务价格，进口品不属于国内生产，因此进口消费品涨价不直接进入 GDP 平减指数。由此，若涨价主要集中在进口消费品，CPI 可能明显上升，而 GDP 平减指数上升较少甚至几乎不变。
\end{solution}
\end{question}

\begin{question}
解释补偿性工资差异与人力资本差异的区别。

\begin{solution}
补偿性工资差异是指不同工作条件、危险程度、舒适度或稳定性不同，企业需要用更高工资补偿较差岗位的非货币劣势；它强调的是\textbf{岗位属性不同}。人力资本差异则强调劳动者教育、培训、经验和技能不同，从而导致生产率与工资差异；它强调的是\textbf{劳动者能力不同}。因此，同样是高工资，一种可能是“因为工作更辛苦更危险”，另一种可能是“因为劳动者更有生产率”。
\end{solution}
\end{question}

\begin{question}
说明在竞争性劳动市场与买方垄断劳动市场中，最低工资对就业影响为什么可能不同。

\begin{solution}
在竞争市场中，工资由供需决定；若最低工资高于均衡工资，会造成劳动供给大于劳动需求，从而出现失业。在买方垄断市场中，企业面临向上倾斜的劳动供给曲线，边际雇佣成本高于工资，因而雇佣量低于竞争水平。此时适度最低工资可能降低边际雇佣成本，使企业愿意雇佣更多工人，就业反而可能上升。关键差异在于：\textbf{竞争市场原本已在均衡点附近，买方垄断市场原本存在雇佣不足。}
\end{solution}
\end{question}

\begin{question}
为什么说基尼系数下降并不必然意味着贫困问题已经解决？

\begin{solution}
基尼系数衡量的是收入分配的不平等程度，是一个相对差距指标；贫困则关注是否低于某个基本生活标准，是绝对生活水平问题。基尼系数下降可能只是高收入群体与中等收入群体差距缩小，而最底层群体收入未必真正提高；也可能是在整体收入都下降时差距缩小，但贫困反而加重。因此降低不平等与减少贫困相关，但并不是同一件事。
\end{solution}
\end{question}

\textbf{三、综合计算题}

\begin{question}
某国 2026 年的宏观数据如下（单位：亿元）：

\begin{itemize}[leftmargin=2em,itemsep=0.2em,topsep=0.2em]
  \item 居民消费 $C=3200$；
  \item 总投资 $I=900$（其中存货投资 100）；
  \item 政府购买 $G=1100$，另有转移支付 200；
  \item 出口 $X=600$，进口 $M=800$；
  \item 折旧（资本消耗）为 250；
  \item 本国居民从国外获得要素收入 120，外国居民从本国获得要素收入 270；
  \item 年末人口为 5000 万人。
\end{itemize}

\begin{enumerate}[label=(\arabic*),leftmargin=2.5em]
  \item 计算 GDP；
  \item 计算 GNP（或 GNI）；
  \item 计算 NDP；
  \item 计算人均 GDP；
  \item 说明为什么转移支付不进入支出法中的 $G$。
\end{enumerate}

\begin{solution}
GDP 由支出法计算：
\[
GDP=C+I+G+X-M=3200+900+1100+600-800=5000.
\]
\ansmath{GDP=5000.}

净国外要素收入为
\[
NFIA=120-270=-150.
\]
因此
\[
GNP=GDP+NFIA=5000-150=4850.
\]
\ansmath{GNP=4850.}

净国内生产总值
\[
NDP=GDP-\text{折旧}=5000-250=4750.
\]
\ansmath{NDP=4750.}

人均 GDP 为
\[
\frac{5000\ \text{亿元}}{5000\ \text{万人}}
=10\ \text{万元/人}.
\]
\ansmath{\text{人均 GDP}=10\ \text{万元/人}.}

转移支付不进入 $G$，因为它不是政府对当期最终产品和服务的购买，不对应当期新增生产；它只是收入再分配，会通过居民消费等渠道间接影响 GDP。
\end{solution}
\end{question}

\begin{question}
某经济体只生产两种商品 $X,Y$。数据如下：

\begin{center}
\begin{tabular}{c|cc}
年份 & $X$（数量，价格） & $Y$（数量，价格）\\ \hline
2025 & $(10,2)$ & $(10,4)$\\
2026 & $(8,3)$ & $(12,5)$
\end{tabular}
\end{center}

\begin{enumerate}[label=(\arabic*),leftmargin=2.5em]
  \item 计算 2025、2026 年名义 GDP；
  \item 以 2025 年为基期，计算 2026 年真实 GDP 与真实增长率；
  \item 计算从 2025 到 2026 的拉氏价格指数、帕氏价格指数与费雪价格指数；
  \item 说明在消费者发生替代行为时，为什么拉氏指数通常高于帕氏指数。
\end{enumerate}

\begin{solution}
名义 GDP：
\[
NGDP_{2025}=10\times 2+10\times 4=60,
\]
\[
NGDP_{2026}=8\times 3+12\times 5=84.
\]
\ansmath{NGDP_{2025}=60,\quad NGDP_{2026}=84.}

以 2025 年为基期，2026 年真实 GDP 为
\[
RGDP_{2026}^{(2025)}=8\times 2+12\times 4=64.
\]
故真实增长率
\[
\frac{64}{60}-1=\frac{4}{60}\approx 6.67\%.
\]
\ansmath{RGDP_{2026}^{(2025)}=64,\quad g_{\text{real}}\approx 6.67\%.}

拉氏价格指数（以 2025 数量作权重）为
\[
P_L=\frac{10\times 3+10\times 5}{10\times 2+10\times 4}
=\frac{80}{60}=1.3333.
\]
帕氏价格指数（以 2026 数量作权重）为
\[
P_P=\frac{8\times 3+12\times 5}{8\times 2+12\times 4}
=\frac{84}{64}=1.3125.
\]
费雪价格指数为
\[
P_F=\sqrt{P_LP_P}
=\sqrt{1.3333\times 1.3125}\approx 1.3229.
\]
故三种通胀率分别约为
\[
33.33\%,\quad 31.25\%,\quad 32.29\%.
\]
\ansmath{P_L\approx 133.33,\quad P_P=131.25,\quad P_F\approx 132.29.}

当消费者发生替代行为时，会减少对涨价较多商品的购买。拉氏指数固定使用基期数量，没有反映这种替代，因此通常会更高；帕氏指数使用当期数量，部分反映了替代行为，所以通常较低。
\end{solution}
\end{question}

\begin{question}
某封闭经济的可贷资金市场由下列函数描述（利率单位为百分数）：
\[
S=120+15r,\qquad I=360-15r.
\]

\begin{enumerate}[label=(\arabic*),leftmargin=2.5em]
  \item 求初始均衡真实利率与均衡投资；
  \item 若政府出现 60 的预算赤字，等价地使国民储蓄函数变为
  \[
  S'=60+15r,
  \]
  求新的均衡真实利率与投资；
  \item 计算预算赤字导致的“挤出效应”大小；
  \item 若此时预期通胀率由 $3\%$ 上升到 $5\%$，且真实利率保持在新的均衡值，名义利率将变为多少？
\end{enumerate}

\begin{solution}
初始均衡由
\[
120+15r=360-15r
\]
得到
\[
30r=240,\qquad r=8.
\]
此时
\[
I=360-15\times 8=240.
\]
\ansmath{r_0=8\%,\quad I_0=240.}

出现预算赤字后，
\[
60+15r=360-15r
\]
得到
\[
30r=300,\qquad r_1=10.
\]
于是
\[
I_1=360-15\times 10=210.
\]
\ansmath{r_1=10\%,\quad I_1=210.}

挤出效应即投资减少量：
\[
240-210=30.
\]
\ansmath{\text{投资被挤出 }30.}

若预期通胀率由 $3\%$ 上升到 $5\%$，真实利率保持在 $10\%$，则按费雪方程
\[
i\approx r+\pi^e=10\%+5\%=15\%.
\]
\ansmath{i\approx 15\%.}
\end{solution}
\end{question}

\begin{question}
考虑一个简化的 AK 增长模型：
\[
Y_t=AK_t,\qquad A=0.20.
\]
资本完全由储蓄形成，折旧率为 $\delta=0.05$，且无人口增长。若储蓄率为 $s$，则资本积累方程为
\[
K_{t+1}-K_t=sY_t-\delta K_t.
\]

\begin{enumerate}[label=(\arabic*),leftmargin=2.5em]
  \item 证明该模型下资本增长率满足 $g_K=sA-\delta$；
  \item 当 $s=0.30$ 时，求长期增长率；
  \item 若通过制度改革把储蓄率提高到 $0.35$，新的长期增长率是多少？
  \item 说明该结果与索罗模型“储蓄率只影响水平、不影响长期增长率”的差别。
\end{enumerate}

\begin{solution}
由
\[
K_{t+1}-K_t=sY_t-\delta K_t=sAK_t-\delta K_t=(sA-\delta)K_t,
\]
两边同时除以 $K_t$，得到
\[
g_K=\frac{K_{t+1}-K_t}{K_t}=sA-\delta.
\]
由于 $Y_t=AK_t$，资本与产出同比例增长，所以
\[
g_Y=g_K=sA-\delta.
\]

当 $s=0.30$ 时，
\[
g=0.30\times 0.20-0.05=0.06-0.05=0.01.
\]
即长期增长率为
\ansmath{g=1\%.}

当 $s=0.35$ 时，
\[
g'=0.35\times 0.20-0.05=0.07-0.05=0.02,
\]
即
\ansmath{g'=2\%.}

与索罗模型不同，AK 模型中资本不存在边际报酬递减，因此更高储蓄率不仅能提高产出水平，还能持续提高长期增长率；索罗模型中由于资本边际报酬递减，更高储蓄率只能提高稳态水平，不能永久提高长期人均增长率（若无外生技术进步）。
\end{solution}
\end{question}

\begin{question}
有两种风险资产 $C$ 与 $D$，两种状态概率都为 $0.5$。收益率如下：

\begin{center}
\begin{tabular}{c|cc}
状态 & 资产 $C$ & 资产 $D$\\ \hline
1 & $40\%$ & $-10\%$\\
2 & $-20\%$ & $20\%$
\end{tabular}
\end{center}

\begin{enumerate}[label=(\arabic*),leftmargin=2.5em]
  \item 分别计算 $C$ 与 $D$ 的期望收益率；
  \item 设投资者将权重 $w$ 配置给 $C$，其余 $1-w$ 配置给 $D$。写出组合在两种状态下的收益率；
  \item 求使组合无风险的权重 $w$ 及此时的无风险收益率；
  \item 解释为什么现实中这种“完美对冲”很少见。
\end{enumerate}

\begin{solution}
资产 $C$ 的期望收益率为
\[
E(r_C)=0.5\times 40\%+0.5\times (-20\%)=10\%.
\]
资产 $D$ 的期望收益率为
\[
E(r_D)=0.5\times (-10\%)+0.5\times 20\%=5\%.
\]
\ansmath{E(r_C)=10\%,\quad E(r_D)=5\%.}

组合在状态 1 的收益率为
\[
r_1=0.4w+(-0.1)(1-w)=0.5w-0.1.
\]
状态 2 的收益率为
\[
r_2=(-0.2)w+0.2(1-w)=0.2-0.4w.
\]

若组合无风险，则两种状态收益率相等：
\[
0.5w-0.1=0.2-0.4w.
\]
解得
\[
0.9w=0.3,\qquad w=\frac13.
\]
此时无风险收益率为
\[
r_f=0.5\times \frac13-0.1=\frac16-0.1=\frac{1}{15}\approx 6.67\%.
\]
\ansmath{w=\frac13,\qquad r_f\approx 6.67\%.}

现实中资产之间极少呈现如此精确的负相关关系；即使在历史数据中看起来能对冲，相关性也会随宏观环境变化而改变，且共同的系统性风险往往无法被完全抵消，所以完美对冲很少真正稳定存在。
\end{solution}
\end{question}

